% GENERATED FILE. Edit canonical lessons, then run npm run generate:books.
% canonical-source-hash: fnv1a64:7a684122566661fb
% canonical-lessons: RU-C195-priroda, RU-C195-vetka, RU-C195-koren, RU-C195-roza, RU-C195-dub

\chapter{Nature, Branch, Root, Rose, Oak}
\label{ch:ru-nature-branch-root-rose-oak}
\hlchaptermodality{195}

\begin{chapteropening}
\textbf{By the end of this chapter:} \emph{I can say nature, a branch, a root, a rose and an oak.}

Complete the last lesson of chapter 195: I can say nature, a branch, a root, a rose and an oak.
\end{chapteropening}

\section[\texorpdfstring{\ru{природа} (priróda)}{природа (priróda)}]{\ru{природа} (priróda) — nature}

\label{lesson:RU-C195-priroda}



\begin{quote}
Before the new one: say the Russian for beetroot, then the Russian for peas.
\end{quote}

\subsection*{You'll want to know: \ru{природа}}
\textbf{\ru{природа}} (\emph{priróda}) — \textquotedblleft{}nature\textquotedblright{}.

\begin{grammarlens}[title={What you've built}]
One more word for this chapter.
\end{grammarlens}

\subsection*{Your turn}
\begin{itemize}
\raggedright
  \item \emph{Say it:} \emph{priróda}
  \item \emph{Say it:} \emph{priróda}, once more
  \item \emph{Say it:} say \emph{priróda}
  \item \emph{Recall:} say the Russian for beetroot, then the Russian for peas, then say \emph{priróda} again
\end{itemize}

\begin{tcolorbox}[breakable,colback=teal!4,colframe=teal!35!black,title={Before you move on}]
What does \ru{природа} mean? (\textquotedblleft{}Nature\textquotedblright{}.) Say it once more.
\end{tcolorbox}
\section[\texorpdfstring{\ru{ветка} (vétka)}{ветка (vétka)}]{\ru{ветка} (vétka) — a branch}

\label{lesson:RU-C195-vetka}



\begin{quote}
Before the new one: say the Russian for peas, then the Russian for nature.
\end{quote}

\subsection*{You'll want to know: \ru{ветка}}
\textbf{\ru{ветка}} (\emph{vétka}) — \textquotedblleft{}a branch\textquotedblright{}.

\begin{grammarlens}[title={What you've built}]
One more word for this chapter.
\end{grammarlens}

\subsection*{Your turn}
\begin{itemize}
\raggedright
  \item \emph{Say it:} \emph{vétka}
  \item \emph{Say it:} \emph{vétka}, once more
  \item \emph{Say it:} say \emph{vétka}
  \item \emph{Recall:} say the Russian for peas, then the Russian for nature, then say \emph{vétka} again
\end{itemize}

\begin{tcolorbox}[breakable,colback=teal!4,colframe=teal!35!black,title={Before you move on}]
What does \ru{ветка} mean? (\textquotedblleft{}A branch\textquotedblright{}.) Say it once more.
\end{tcolorbox}
\section[\texorpdfstring{\ru{корень} (kóren)}{корень (kóren)}]{\ru{корень} (kóren) — a root}

\label{lesson:RU-C195-koren}



\begin{quote}
Before the new one: say the Russian for nature, then the Russian for a branch.
\end{quote}

\subsection*{You'll want to know: \ru{корень}}
\textbf{\ru{корень}} (\emph{kóren}) — \textquotedblleft{}a root\textquotedblright{}.

\begin{grammarlens}[title={What you've built}]
One more word for this chapter.
\end{grammarlens}

\subsection*{Your turn}
\begin{itemize}
\raggedright
  \item \emph{Say it:} \emph{kóren}
  \item \emph{Say it:} \emph{kóren}, once more
  \item \emph{Say it:} say \emph{kóren}
  \item \emph{Recall:} say the Russian for nature, then the Russian for a branch, then say \emph{kóren} again
\end{itemize}

\begin{tcolorbox}[breakable,colback=teal!4,colframe=teal!35!black,title={Before you move on}]
What does \ru{корень} mean? (\textquotedblleft{}A root\textquotedblright{}.) Say it once more.
\end{tcolorbox}
\section[\texorpdfstring{\ru{роза} (róza)}{роза (róza)}]{\ru{роза} (róza) — a rose}

\label{lesson:RU-C195-roza}



\begin{quote}
Before the new one: say the Russian for a branch, then the Russian for a root.
\end{quote}

\subsection*{You'll want to know: \ru{роза}}
\textbf{\ru{роза}} (\emph{róza}) — \textquotedblleft{}a rose\textquotedblright{}.

\begin{grammarlens}[title={What you've built}]
One more word for this chapter.
\end{grammarlens}

\subsection*{Your turn}
\begin{itemize}
\raggedright
  \item \emph{Say it:} \emph{róza}
  \item \emph{Say it:} \emph{róza}, once more
  \item \emph{Say it:} say \emph{róza}
  \item \emph{Recall:} say the Russian for a branch, then the Russian for a root, then say \emph{róza} again
\end{itemize}

\begin{tcolorbox}[breakable,colback=teal!4,colframe=teal!35!black,title={Before you move on}]
What does \ru{роза} mean? (\textquotedblleft{}A rose\textquotedblright{}.) Say it once more.
\end{tcolorbox}
\section[\texorpdfstring{\ru{дуб} (dub)}{дуб (dub)}]{\ru{дуб} (dub) — an oak}

\label{lesson:RU-C195-dub}



\begin{quote}
Before the new one: say the Russian for a root, then the Russian for a rose.
\end{quote}

\subsection*{You'll want to know: \ru{дуб}}
\textbf{\ru{дуб}} (\emph{dub}) — \textquotedblleft{}an oak\textquotedblright{}.

\begin{grammarlens}[title={What you've built}]
One more word for this chapter.
\end{grammarlens}

\subsection*{Your turn}
\begin{itemize}
\raggedright
  \item \emph{Say it:} \emph{dub}
  \item \emph{Say it:} \emph{dub}, once more
  \item \emph{Say it:} say \emph{dub}
  \item \emph{Recall:} say the Russian for a root, then the Russian for a rose, then say \emph{dub} again
\end{itemize}

\begin{tcolorbox}[breakable,colback=teal!4,colframe=teal!35!black,title={Before you move on}]
What does \ru{дуб} mean? (\textquotedblleft{}An oak\textquotedblright{}.) Say it once more.
\end{tcolorbox}
